\documentclass[12pt]{article}
\usepackage[utf8]{inputenc}
\usepackage[T1]{fontenc}
\usepackage{hyperref}
\title{Bajjika Language Resources for K-12 Curriculum Development}
\author{Compiled by Claude Sonnet 5 (claude-sonnet-5) for Andr\'as Kornai}
\date{}
\begin{document}
\maketitle

\section{Introduction}
Bajjika (ISO 639-3 vjk; also Tirhutiya) is spoken mainly in north-western Bihar, India (Samastipur, Sitamarhi, Muzaffarpur, Vaishali, Sheohar) and in Nepal. Nepal's 2021 census counted about 1.13 million speakers there, per Wikipedia. Sources disagree on whether it is a language or a dialect of Maithili. Documentation is thin, and the search found no school curriculum.

\section{Primary Educational Resources}
None found in the search results.

\section{Linguistic and Reference Materials}
\begin{itemize}
\item Wikipedia, Bajjika: \url{https://en.wikipedia.org/wiki/Bajjika}
\item Wikipedia, Languages of Bihar: \url{https://en.wikipedia.org/wiki/Languages_of_Bihar}
\item Bharati Bajjika Hindi Kos (2008), edited by Awadheshwar Arun and Ram Niwas Sharma, a Bajjika--Hindi dictionary. It appears in a reference list at: \url{https://lwc1.benjamins.com/catalog/pbns.264.07kum}
\item Abhishek Kumar Kashyap's 2009 PhD dissertation on Bajjika, and his articles on Bajjika, appear in the same reference list.
\item George Grierson, \emph{Seven Grammars of the Dialects and Subdialects of the Bihari Language} (1883--84), a historical work: \url{https://catalogue.bnf.fr/ark:/12148/cb30541389q}
\end{itemize}

\section{Religious and Community Resources}
None found in the search results.

\section{Digital and Audio Resources}
\begin{itemize}
\item Bajjika: Swadesh List Elicitation Sessions, an ELDP archive collection held at SOAS: \url{https://eldp.access.preservica.com/dk0038/}
\end{itemize}

\section{Curriculum Development Context}
Bajjika has many speakers but little published documentation or teaching material. Instruction in Bihar normally uses Hindi, and Bajjika is often treated as a Maithili dialect. A curriculum would need new work on orthography and teaching materials, and a decision on script and relation to Hindi and Maithili.
\end{document}